Emmy wist niet waarom,
maar ze keek steeds rondom.

“Wie snijdt er met grijnslach
op deze niet-zondag

door mijn ijzige vlees?”
was haar vragende vrees.

Een afgescheidenheid
van haar zelfbewustheid.

Onze basale angst.
Een ongewenste vangst.

Gedreven door instinct,
dat haar gluurder verlinkt.

Als ze werd bekeken
dan kon ze uitsteken

haar scherpe voelhorens.
Zo wist ze alvorens

het bijzijn van een vent.
Dat gebeurde frequent.

Oud of jong, goor of knap,
wierp om de haverklap

haar stiekem een blik toe.
Blijkbaar was dat taboe.

Als ze werden gesnapt,
dan hapten ze betrapt.

Via dit vermogen
ervoer ze blikogen,

over haar naakte lijf.
Ze staakte haar verblijf

samen met eendenkroos,
dat plakte op haar doos.

Ze keek dit keer wat pips.
Wiebelend met haar bips

verliet ze de rietven
huppelend naar een den,

waar haar vodden hingen
en andere dingen.

Licht over het veenmos
naar het grauwe veenbos.

Haar tietjes deinden mee.
Haar lokken ook gedwee.

Het was voorheen een kerk.
Daarna deed het zijn werk

als een kenniscentrum
met als curriculum

het lezen en schrijven.
Het geloof bedrijven

en het hoofdrekenen
kon men op rekenen.

Werd men zo bekwamer?
Haar school had één kamer.

Het strakke onderricht
werd nauwelijks verlicht

door vensters naar buiten.
Meisjes sloot men buiten.

Niet Emmy en Emma.
Dat bleek een dilemma

voor de rest van het dorp,
dat lag op een steenworp

afstand van de zuidpool
van deze kloosterschool.

Wie was de weldoener?
Hoe werd het gras groener

voor Emmy haar gezin?
Een begijn greep toen in:

“Dat is God’s genade.
Jaloezie het kwade!”

Het was stampend leren
en soms modelleren.

Het begrip bleef achter.
De leraar werd wachter.

Hij pushte prestatie
middels competitie.

Strenge discipline
met doodse routine.

Hardop leerde men zo
in zijn eigen tempo.

De oudere steunde,
de jongere leunde.

Een schema naar seizoen.
In de winter veel doen

en in de zomer ging
bijna elke leerling

en stak uit zijn werkhand
op het vruchtbare land.

De les was winterkoud
en sneed zelden brandhout.

Het duurde vaak uren.
Men moest het verduren.

Emma gaf het snel op.
Emmy vond de school top

en haar starre rooster
als pad naar het klooster.

“Hé, ben jij ook zwanger?”
Voorts zei de kerkganger:

“Nee! Ik ben gewoon dik!”
Emmy’s moeder had schik.

Haar man bleek zo verheugd
met moeders christendeugd

dat hij het alom zag.
Zijn enige beklag

was hun intimiteit,
die was hij even kwijt.

Ze hadden, pas getrouwd,
net hun band opgebouwd,

toen moeders zwangerschap
hem vroeg om rekenschap.

Hun huwelijkscadeau
suste zijn libido.

De landheer schonk vader
een gratis ontlader,

die zijn lusten stilde,
zo vaak als hij wilde.

Moeder vond het niet erg.
Vermomd als een herberg

was het landheers bordeel.
Hij kwam er destijds veel.

Er was een bekende
hoopvolle legende

in Emmy’s regio,
als een scenario

voor een vogelvrije.
En wel een gastvrije.

Reizigers verhaalden,
hoe zij die ronddwaalden

een vrijhaven vonden,
waar ze schuilen konden,

zonder angst voor wraking,
vervolging of kwelling.

Er huisde een sekte,
die haar hand uitstrekte,

zelfs naar criminelen.
Het wierf daarmee velen.

Het was Emmy’s moeder,
die haar als opvoeder

op de gevaren wees.
“Emmy, waar ik voor vrees,

is dat je dit gelooft.
Denk niet, als een leeghoofd,

dat zo’n plek heus bestaat.
Geen stad is zonder kwaad,

slechts bewoond door braven.
Het is een vrij-haven,

zoals de naam al zegt.
Dit lusthof is niet echt.

Het komt voort uit een boek.
De rest is lariekoek.”

“Haal jij Emma nu op?"
Vader kreeg op zijn kop

omdat hij werd verblind
door alles naast zijn kind.

“Verstrooide professor!”
noemde de spectator

de huisvoogds distractie.
De justificatie:

“Mijn geestigvermogen
moet de mens niet pogen

te onderwaarderen!”
was om te proberen

met wat dure woorden
zijn aard te vermoorden.

Haar lachen niet gesnapt,
want hij had niet gegrapt,

grijnsde hij maar wat mee
en liep langs het cafe

op weg naar het klooster.
Paste in zijn rooster

om een pint te drinken?
Wie zou hem verlinken?

Een rode neus later
ging hij met zijn kater

weer op weg naar zijn huis.
Hij was al even thuis.

Toen schreeuwde Emmy’s ma
“Waar hang je uit, Emma?”

Moeder zag aan zijn blik
de al bekende schrik.

Zijn fout direct gesnapt
snelde vader betrapt

weer op weg naar Emma.
Vader kwam met sperma

vlekken op zijn schaamlap,
thuis na weer een uitstap.

Dit keer naar het bordeel
als bekend ritueel.
Hij hoorde: “En Emma?”,
en gaf een fout lemma:

“Het was slechts absentie
door kontinhetplaatsie!”

Emmy begreep niet goed
haar moeders zwaar gemoed,

wanneer men iets besprak.
Haar verleden doorbrak

bij het naamwoord ‘profeet’
en ma hem eer aan deed

als ze hem citeerde.
Zij blijkbaar ooit leerde:

“Het menselijk lichaam
treft geen enkele blaam,

maar volgt slechts de wetten,
die tijd en al zetten.

Als voertuig op een reis,
door dit aards paradijs,

verlangt, lijdt en sterft het,
na te zijn afgezet.

Het is slechts de tempel,
waar de ziel zijn stempel

goddelijk op afdrukt.
De geest is een product

van het brein als orgaan.
Waar gedachten rondgaan,

sensaties binnengaan
en emoties ontstaan.

De geest is ook een brug
naar de ziel en terug,

afdalend naar het lijf
als zijn binnenverblijf.

Geen lichaam staat ooit los
van God of de kosmos.

Het hoort primair erbij.
Een basale bladzij

van ons grenzeloos boek.
Ook al raakt God soms zoek,

een vonk geheten ziel
vindt altijd een ventiel

om ons hart te vullen.
Aldaar te onthullen

versterkt als door een lens,
zijn liefde zo intens,

dat alles weer verbindt.
Ik hoop dat jij God vindt!”

Neoplatonisme
dualisme, monisme,

panentheïsme met
een gedegen natuurwet

werden door haar profeet
ooit aan elkaar gesmeed.

Emmy was verwonderd,
hoe ma afgezonderd,

die wijsheden ontving,
schoon ze nooit op pad ging.
